\section{Circuit Complexity --- PSubset PPoly Program}
%%% generated by the blueprint pipeline from TCSlib.Complexity.CircuitComplexity.PSubsetPPolyProgram
\subsection*{Circuit Complexity --- PSubset PPoly Program: Introduction}

Declarations of the Lean module \texttt{TCSlib.Complexity.CircuitComplexity.PSubsetPPolyProgram}.

\subsection*{Circuit Complexity --- PSubset PPoly Program: Results}

\begin{definition}[Bit source]
\lean{BoolCircuit.BitSrc}
\label{BoolCircuit.BitSrc}
\leanok
A bit source is one of three kinds: an input bit indexed by a natural number, a Boolean
constant, or bit $j$ of the block computed by instruction $i$, for natural numbers
$i,j$.
\end{definition}

\begin{definition}[Finite-function instruction]
\lean{BoolCircuit.GInstr}
\label{BoolCircuit.GInstr}
\leanok
\uses{BoolCircuit.BitSrc}
For a type $\kappa$, an instruction consists of a kind in $\kappa$ selecting the finite
function to compute and an ordered list of source bits.
\end{definition}

\begin{definition}[Evaluation of a bit source]
\lean{BoolCircuit.BitSrc.eval}
\label{BoolCircuit.BitSrc.eval}
\leanok
\uses{BoolCircuit.BitSrc}
Given an input list $x$ of Boolean values and a list $b$ of Boolean blocks, evaluate a
bit source as follows: input index $k$ reads $x_k$, using false if unavailable; a
constant source returns its Boolean value; and block source $(i,j)$ reads bit $j$ of
block $i$, using an empty block or false bit when the requested entry is unavailable.
\end{definition}

\begin{definition}[Validity of a Bit Source]
\lean{BoolCircuit.BitSrc.Valid}
\label{BoolCircuit.BitSrc.Valid}
\leanok
\uses{BoolCircuit.BitSrc}
For natural numbers $n,w,i$, a bit source is valid if it is an input bit with index less
than $n$, a Boolean constant, or a bit with index $j<w$ from a block computed by an
instruction with index less than $i$.
\end{definition}

\begin{definition}[Running a Boolean Program]
\lean{BoolCircuit.progRun}
\label{BoolCircuit.progRun}
\leanok
\uses{BoolCircuit.BitSrc.eval, BoolCircuit.GInstr}
Given a function $F$ that assigns a list of output bits to each instruction kind and
list of input bits, an input list $x$, a list of instructions, and an initial list of
bit blocks, process the instructions in order. At each instruction, evaluate its ordered
source list using $x$ and the blocks computed so far, apply $F$ with that instruction's
kind, and append the resulting block to the list; the resulting list is the program run.
\end{definition}

\begin{definition}[Blocks Computed by a Program]
\lean{BoolCircuit.progBlocks}
\label{BoolCircuit.progBlocks}
\leanok
\uses{BoolCircuit.GInstr, BoolCircuit.progRun}
Given a function $F$ assigning an output list to each instruction kind and Boolean input
list, and an input list $x$, the blocks computed by a program are the blocks produced by
running its instructions in order, starting with no initial blocks.
\end{definition}

\begin{theorem}[Concatenation law for program execution]
\lean{BoolCircuit.progRun_append}
\label{BoolCircuit.progRun_append}
\leanok
\uses{BoolCircuit.GInstr, BoolCircuit.progRun}
\difficulty{1}
For any instruction-kind type, function $F$ assigning a Boolean-list function to each
kind, input list $x$, instruction lists $p,q$, and initial list of Boolean blocks,
executing the concatenated program $p$ followed by $q$ yields the same list of blocks as
first executing $p$ and then executing $q$ from the resulting state.
\end{theorem}

\begin{theorem}[Length of the program run]
\lean{BoolCircuit.length_progRun}
\label{BoolCircuit.length_progRun}
\leanok
\uses{BoolCircuit.BitSrc.eval, BoolCircuit.GInstr, BoolCircuit.progRun}
\difficulty{4}
Let $F$ assign to each instruction kind a function from Boolean lists to Boolean lists.
For every input list $x$, program $prog$, and initial list of Boolean blocks $init$,
running the program appends one block per instruction, so the resulting list of blocks
has length $|init|+|prog|$.
\end{theorem}

\begin{theorem}[Preservation of initial blocks]
\lean{BoolCircuit.progRun_getD_of_lt}
\label{BoolCircuit.progRun_getD_of_lt}
\leanok
\uses{BoolCircuit.BitSrc.eval, BoolCircuit.GInstr, BoolCircuit.progRun}
\difficulty{4}
For any instruction-kind type, function $F$ assigning a Boolean-list function to each
kind, input list $x$, program, initial list of Boolean blocks, and index $i$ less than
the length of the initial list, the block at index $i$ after executing the program
equals the initial block at that index.
\end{theorem}

\begin{definition}[Well-Formed Boolean Program]
\lean{BoolCircuit.ProgWF}
\label{BoolCircuit.ProgWF}
\leanok
\uses{BoolCircuit.BitSrc.Valid, BoolCircuit.GInstr}
Let $n,w$ be natural numbers, let $F$ assign an output list to each instruction kind and
Boolean input list, and let a program be a finite list of instructions. The program is
well formed if, for every instruction at position $i$ in the program and every source
used by that instruction, each input-bit index is less than $n$, every constant is
allowed, and every referenced earlier block has index less than $i$ and bit index less
than $w$.
\end{definition}

\begin{theorem}[Evaluation of a valid bit source]
\lean{BoolCircuit.BitSrc.eval_congr}
\label{BoolCircuit.BitSrc.eval_congr}
\leanok
\uses{BoolCircuit.BitSrc, BoolCircuit.BitSrc.Valid, BoolCircuit.BitSrc.eval}
\difficulty{3}
Let $x$ be a finite Boolean list, let $n,w,i$ be natural numbers, and let $s$ be a bit
source valid for input length $n$, block width $w$, and instruction index $i$. If two
lists of previously computed Boolean blocks agree at every index $i'<i$, then evaluating
$s$ using either list gives the same Boolean value.
\end{theorem}

\begin{theorem}[Computed block at a program position]
\lean{BoolCircuit.progBlocks_getD}
\label{BoolCircuit.progBlocks_getD}
\leanok
\uses{BoolCircuit.BitSrc.eval, BoolCircuit.BitSrc.eval_congr, BoolCircuit.GInstr, BoolCircuit.ProgWF, BoolCircuit.progBlocks, BoolCircuit.progRun, BoolCircuit.progRun_append, BoolCircuit.progRun_getD_of_lt}
\difficulty{5}
Let $F$ map an instruction kind and a Boolean list to a Boolean list, let $x$ be a
finite Boolean list, and let $\mathrm{prog}$ be a finite list of instructions. Suppose
the program is well formed for input bound $n$ and block width $w$: each source used at
position $i$ is an input position below $n$, a constant, or a bit below $w$ in an
earlier block. At every valid instruction position $i$, the block computed there is $F$
applied to that instruction's kind and the list obtained by evaluating its sources using
$x$ and the program's computed blocks.
\end{theorem}

\begin{definition}[Gadgets for an Instruction Kind]
\lean{BoolCircuit.kindGadgets}
\label{BoolCircuit.kindGadgets}
\leanok
\uses{BoolCircuit.DAGCircuit, BoolCircuit.gadget}
Let $m,w$ be natural numbers, let $F$ assign an output list to each instruction kind and
Boolean input list, and fix an instruction kind $\kappa_0$. The associated list
consists, for each $j<w$, of the Boolean circuit gadget for the function that sends an
input vector $v$ of length $m$ to the bit at position $j$ of $F(\kappa_0,v)$, taking
that bit to be false when the output list has no position $j$.
\end{definition}

\begin{theorem}[Number of kind gadgets]
\lean{BoolCircuit.length_kindGadgets}
\label{BoolCircuit.length_kindGadgets}
\leanok
\uses{BoolCircuit.kindGadgets}
\difficulty{1}
Let $\kappa$ be a type, let $m,w\in\mathbb{N}$, let $F$ assign a Boolean list to each
pair consisting of an element of $\kappa$ and a Boolean list, and choose
$\kappa_0\in\kappa$. The list of gadgets associated with $\kappa_0$ has length $w$.
\end{theorem}

\begin{definition}[Width associated with a kind]
\lean{BoolCircuit.kindWidth}
\label{BoolCircuit.kindWidth}
\leanok
\uses{BoolCircuit.embedWidth, BoolCircuit.kindGadgets}
Let $\kappa$ be a type, let $m$ and $w$ be natural numbers, and let $F$ assign a Boolean
list to each pair consisting of an element of $\kappa$ and a Boolean list. For each
$\kappa_0\in\kappa$, define its kind width to be the embedding width of the kind gadgets
determined by $m$, $w$, $F$, and $\kappa_0$.
\end{definition}

\begin{definition}[Base Vertex of a Program Instruction]
\lean{BoolCircuit.progBase}
\label{BoolCircuit.progBase}
\leanok
\uses{BoolCircuit.GInstr, BoolCircuit.kindWidth}
Let $m,w,n$ be natural numbers, let $F$ assign an output list to each instruction kind
and Boolean input list, and let a program be a finite list of instructions. The base
vertex at position $i$ is $n+2$ plus the sum of the widths associated with the kinds of
the first $i$ instructions (or all available instructions if the program has fewer than
$i$).
\end{definition}

\begin{theorem}[Monotonicity of program bases]
\lean{BoolCircuit.progBase_mono}
\label{BoolCircuit.progBase_mono}
\leanok
\uses{BoolCircuit.GInstr, BoolCircuit.kindWidth, BoolCircuit.progBase}
\difficulty{4}
Let $n,m,w$ be nonnegative integers, let $\kappa$ be a type of instruction kinds, let
$F$ assign an output list of Booleans to each kind and input list, and let
$\mathrm{prog}$ be a finite list of instructions with kinds in $\kappa$. For each kind,
let its width be the total number of gates in the associated gadgets. The base offset
$n+2$ plus the sum of the widths of the first $i$ instructions is nondecreasing in $i$:
if $i\le i'$, then the base at $i$ is at most the base at $i'$.
\end{theorem}

\begin{theorem}[Successive program base offsets]
\lean{BoolCircuit.progBase_succ}
\label{BoolCircuit.progBase_succ}
\leanok
\uses{BoolCircuit.GInstr, BoolCircuit.kindWidth, BoolCircuit.progBase}
\difficulty{3}
Let $n,m,w$ be nonnegative integers, let $\kappa$ be a type of instruction kinds, let
$F$ assign an output list of Booleans to each kind and input list, and let
$\mathrm{prog}$ be a finite list of instructions with kinds in $\kappa$. For each kind,
let its width be the total number of gates in the associated gadgets. If $i$ is a valid
position in the program, then the base offset for the first $i+1$ instructions equals
the base offset for the first $i$ instructions plus the width of the kind of instruction
$i$.
\end{theorem}

\begin{definition}[Vertex for a Program Output Bit]
\lean{BoolCircuit.progVertex}
\label{BoolCircuit.progVertex}
\leanok
\uses{BoolCircuit.GInstr, BoolCircuit.embedOffset, BoolCircuit.kindGadgets, BoolCircuit.progBase}
Let $m,w,n$ be natural numbers, let $F$ assign an output list to each instruction kind
and Boolean input list, let a program be a finite list of instructions, and assume the
type of instruction kinds has a default element. The vertex associated with instruction
position $i$ and bit position $j$ is the base vertex at $i$ plus the embedding offset at
$j$ in the gadget list for the kind of instruction $i$, using the default instruction if
position $i$ is beyond the program's end.
\end{definition}

\begin{definition}[Vertex assigned to a source bit]
\lean{BoolCircuit.srcVertex}
\label{BoolCircuit.srcVertex}
\leanok
\uses{BoolCircuit.BitSrc, BoolCircuit.BitSrc.Valid, BoolCircuit.GInstr, BoolCircuit.progVertex}
Let $n,m,w$ be natural numbers, let $\kappa$ be a type with a chosen element, let $F$
assign an output bit list to each element of $\kappa$ and input bit list, and let
$\mathrm{prog}$ be a finite list of instructions. For instruction index $i$ and source $s$, the
assigned vertex is its input index when $s$ is a valid input source, is $n$ or $n+1$ for
the valid constant false or true source, respectively, and is the vertex assigned to the
referenced earlier instruction and output bit when $s$ is a valid block source; every
invalid source is assigned vertex $n$.
\end{definition}

\begin{definition}[Source vertices for an instruction]
\lean{BoolCircuit.srcVertices}
\label{BoolCircuit.srcVertices}
\leanok
\uses{BoolCircuit.GInstr, BoolCircuit.srcVertex}
Let $n,m,w$ be natural numbers, let $\kappa$ be a type with a chosen element, let $F$
assign an output bit list to each element of $\kappa$ and input bit list, and let
$\mathrm{prog}$ be a finite instruction list. For instruction index $i$, form the list of
vertices assigned to its sources, append $m$ copies of the vertex $n$, and retain the
first $m$ entries; the resulting list is the list of source vertices for that
instruction.
\end{definition}

\begin{theorem}[Length of the source-vertex list]
\lean{BoolCircuit.length_srcVertices}
\label{BoolCircuit.length_srcVertices}
\leanok
\uses{BoolCircuit.GInstr, BoolCircuit.srcVertices}
\difficulty{1}
Let $m,w,n$ be natural numbers, let $F$ assign to each instruction kind a function from
Boolean lists to Boolean lists, and let $prog$ be a program of such instructions. For
every instruction index $i$, the list of source vertices associated with that
instruction has length $m$.
\end{theorem}

\begin{definition}[Gate list for a program prefix]
\lean{BoolCircuit.progGates}
\label{BoolCircuit.progGates}
\leanok
\uses{BoolCircuit.DAGGate, BoolCircuit.GInstr, BoolCircuit.constGate, BoolCircuit.embedAll, BoolCircuit.kindGadgets, BoolCircuit.progBase, BoolCircuit.srcVertices}
Let $n,m,w$ be natural numbers, let $\kappa$ be a type with a chosen element, let $F$
assign a finite Boolean list to each pair consisting of an element of $\kappa$ and a
finite Boolean list, and let $\mathrm{prog}$ be a finite list of instructions, each
carrying an element of $\kappa$. Define a gate list $G_i$ for every natural number $i$
by taking $G_0$ to consist of the constant-false and constant-true gates, and, for each
$i$, forming $G_{i+1}$ by appending to $G_i$ the embedded gadgets for the kind carried
by instruction $i$ (using the chosen element when that instruction is absent), with
source vertices determined by that instruction and base offset $n+2$ plus the total
widths of the preceding instructions.
\end{definition}

\begin{theorem}[Gate count and program offset]
\lean{BoolCircuit.length_progGates}
\label{BoolCircuit.length_progGates}
\leanok
\uses{BoolCircuit.GInstr, BoolCircuit.kindWidth, BoolCircuit.length_embedAll, BoolCircuit.progBase, BoolCircuit.progGates}
\difficulty{4}
Let $m,w,n$ be natural numbers, let $\kappa$ be a type with a chosen default element, let $F$
assign a Boolean-list function to each element of $\kappa$, and let $P$ be a finite list of
instructions whose kinds lie in $\kappa$. For each index $i$ with $i<|P|$, the number of
input vertices plus the number of gates in the circuit generated from the first $i$
instructions equals the base vertex index assigned to stage $i$.
\end{theorem}

\begin{theorem}[Program output vertex bound]
\lean{BoolCircuit.progVertex_lt}
\label{BoolCircuit.progVertex_lt}
\leanok
\uses{BoolCircuit.GInstr, BoolCircuit.embedOffset_lt, BoolCircuit.kindGadgets, BoolCircuit.kindWidth, BoolCircuit.progBase, BoolCircuit.progBase_succ, BoolCircuit.progVertex}
\difficulty{3}
Let $n,m,w$ be natural numbers, let the instruction kinds form a type with a designated
default, and let $F$ assign a Boolean-list function to each kind. For any program,
instruction index $i$ less than the program length, and output-bit index $j<w$, the
vertex assigned to bit $j$ of instruction $i$ is strictly less than the base vertex
immediately after that instruction.
\end{theorem}

\begin{theorem}[Source vertices precede the instruction block]
\lean{BoolCircuit.srcVertex_lt}
\label{BoolCircuit.srcVertex_lt}
\leanok
\uses{BoolCircuit.BitSrc, BoolCircuit.BitSrc.Valid, BoolCircuit.GInstr, BoolCircuit.progBase, BoolCircuit.progBase_mono, BoolCircuit.progVertex_lt, BoolCircuit.srcVertex}
\difficulty{5}
Let $K$ be a nonempty set of instruction kinds, let $m,w,n$ be natural numbers, and let
$F:K\to\{0,1\}^*\to\{0,1\}^*$ assign an output bit string to each kind and input bit
string. A program is a finite list of instructions, each with a kind in $K$ and an
ordered list of bit sources. For a kind $k$, let $W(k)$ be the sum, over the first $w$
output positions, of one plus the number of gates in the circuit compiled from the
Boolean function giving that output bit of $F(k,\cdot)$ on $m$ input bits (an
out-of-range output position is read as zero). For an index $i$ no greater than the
program length, set $B_i=n+2+\sum_{r<i}W(k_r)$, where $k_r$ is the kind of instruction
$r$. For every bit source $s$, its assigned vertex is strictly less than $B_i$: a valid
input source $k$ is assigned vertex $k$, a valid constant is assigned vertex $n$ or
$n+1$ according as it is false or true, and a valid reference to output position $j$ of
an earlier instruction $i'$ is assigned vertex $B_{i'}$ plus that position's offset in
the corresponding output circuits; an invalid source is assigned vertex $n$.
\end{theorem}

\begin{theorem}[Source list vertices precede the instruction block]
\lean{BoolCircuit.srcVertices_lt}
\label{BoolCircuit.srcVertices_lt}
\leanok
\uses{BoolCircuit.GInstr, BoolCircuit.progBase, BoolCircuit.srcVertex_lt, BoolCircuit.srcVertices}
\difficulty{4}
Let $K$ be a nonempty set of instruction kinds, let $m,w,n$ be natural numbers, and let
$F:K\to\{0,1\}^*\to\{0,1\}^*$ assign an output bit string to each kind and input bit
string. A program is a finite list of instructions, each with a kind in $K$ and an
ordered list of bit sources. For each kind $k$, let $W(k)$ be the sum, over the first
$w$ output positions, of one plus the number of gates in the circuit compiled from the
Boolean function giving that output bit of $F(k,\cdot)$ on $m$ input bits (an
out-of-range output position is read as zero). For $i$ no greater than the program
length, put $B_i=n+2+\sum_{r<i}W(k_r)$, where $k_r$ is the kind of instruction $r$. Form
a list by mapping each source of instruction $i$ to its assigned vertex, appending $m$
copies of $n$, and retaining only the first $m$ entries. Every vertex in this list is
strictly less than $B_i$.
\end{theorem}

\begin{theorem}[Acyclicity of program gates]
\lean{BoolCircuit.gatesAcyclic_progGates}
\label{BoolCircuit.gatesAcyclic_progGates}
\leanok
\uses{BoolCircuit.GInstr, BoolCircuit.GatesAcyclic, BoolCircuit.GatesAcyclic.append, BoolCircuit.constGate, BoolCircuit.gatesAcyclic_embedAll, BoolCircuit.length_progGates, BoolCircuit.length_srcVertices, BoolCircuit.progGates, BoolCircuit.srcVertices_lt}
\difficulty{4}
Let $n$ be the number of input bits, let $\mathrm{prog}$ be a finite list of
instructions, and let $i\in\mathbb{N}$ satisfy $i\leq |\mathrm{prog}|$. The gate list
produced for the first $i$ instruction stages is acyclic relative to the $n$ input
vertices.
\end{theorem}

\begin{theorem}[Fan-in of program gates]
\lean{BoolCircuit.faninTwo_progGates}
\label{BoolCircuit.faninTwo_progGates}
\leanok
\uses{BoolCircuit.GInstr, BoolCircuit.constGate, BoolCircuit.faninTwo_embedAll, BoolCircuit.gadget_isFaninTwo, BoolCircuit.kindGadgets, BoolCircuit.progGates}
\difficulty{4}
Let $n$ be a natural number, let $\kappa$ be a type of instruction kinds, and let
$\mathrm{prog}$ be a finite program of instructions of those kinds. For every natural
number $i$, each gate in the circuit constructed from the first $i$ program steps has
distinct input vertices, every NOT gate has exactly one input, and every gate has at
most two inputs.
\end{theorem}

\begin{theorem}[Values computed by program gates]
\lean{BoolCircuit.progGates_values}
\label{BoolCircuit.progGates_values}
\leanok
\uses{BoolCircuit.BitSrc.Valid, BoolCircuit.BitSrc.eval, BoolCircuit.DAGGate.eval, BoolCircuit.GInstr, BoolCircuit.ProgWF, BoolCircuit.embedAll, BoolCircuit.gadget_eval, BoolCircuit.kindGadgets, BoolCircuit.length_progGates, BoolCircuit.length_runWith, BoolCircuit.length_srcVertices, BoolCircuit.ofFn_getD_eq_map, BoolCircuit.progBase, BoolCircuit.progBase_mono, BoolCircuit.progBlocks, BoolCircuit.progBlocks_getD, BoolCircuit.progGates, BoolCircuit.progVertex, BoolCircuit.progVertex_lt, BoolCircuit.runWith, BoolCircuit.runWith_append, BoolCircuit.runWith_embedAll_getD, BoolCircuit.runWith_eval_constGates, BoolCircuit.runWith_getD_of_lt, BoolCircuit.srcVertex, BoolCircuit.srcVertices, BoolCircuit.srcVertices_lt}
\difficulty{6}
Let $n,m,w$ be natural numbers, let the instruction kinds form a type with a designated
default, and let $F$ assign a Boolean-list function to each kind. For an input
assignment $x\in\{\mathrm{false},\mathrm{true}\}^n$, a program of instructions whose
source bits are valid for their instruction positions and whose source lists each have
length $m$, and any $i\leq$ the program length, running the first $i$ stages of the
constructed gate list on $x$ preserves every input bit, places false and true at
vertices $n$ and $n+1$, respectively, and gives at each earlier instruction's output
vertex the corresponding bit of that instruction's computed block.
\end{theorem}

\begin{definition}[Circuit of a program]
\lean{BoolCircuit.progCircuit}
\label{BoolCircuit.progCircuit}
\leanok
\uses{BoolCircuit.DAGCircuit, BoolCircuit.GInstr, BoolCircuit.gatesAcyclic_progGates, BoolCircuit.length_progGates, BoolCircuit.progBase, BoolCircuit.progBase_mono, BoolCircuit.progGates, BoolCircuit.progVertex, BoolCircuit.progVertex_lt}
Let $n,m,w$ be natural numbers, let $\kappa$ be a type with a chosen element, let $F$
assign an output bit list to each element of $\kappa$ and input bit list, and let
$\mathrm{prog}$ be a finite instruction list. For indices $i,j$, the program circuit has the
gate list obtained by constructing the program gates through the full program length.
Its output is the vertex assigned to instruction $i$ and bit $j$ when $i$ is within the
program and $j<w$, and is the input vertex $n$ otherwise.
\end{definition}

\begin{theorem}[Evaluation of a compiled program circuit]
\lean{BoolCircuit.progCircuit_eval}
\label{BoolCircuit.progCircuit_eval}
\leanok
\uses{BoolCircuit.DAGCircuit.eval, BoolCircuit.DAGCircuit.values, BoolCircuit.GInstr, BoolCircuit.ProgWF, BoolCircuit.progBlocks, BoolCircuit.progCircuit, BoolCircuit.progGates_values}
\difficulty{1}
Let $n,m,w$ be nonnegative integers, let $\kappa$ be an inhabited type of instruction
kinds, let $F$ map a kind and a Boolean list to a Boolean list, and let $\mathrm{prog}$
be a finite instruction list. Suppose the program is well formed for input bound $n$ and
block width $w$, and every instruction has exactly $m$ sources. For valid positions
$i<|\mathrm{prog}|$ and $j<w$, and every Boolean assignment $x$ to the $n$ inputs, the
circuit associated with position $i$ and output bit $j$ evaluates to bit $j$ of the
block computed by running the program on $x$, with missing list entries interpreted as
false.
\end{theorem}

\begin{theorem}[Fan-in of the program circuit]
\lean{BoolCircuit.progCircuit_isFaninTwo}
\label{BoolCircuit.progCircuit_isFaninTwo}
\leanok
\uses{BoolCircuit.DAGCircuit.IsFaninTwo, BoolCircuit.GInstr, BoolCircuit.faninTwo_progGates, BoolCircuit.progCircuit}
\difficulty{1}
For any natural numbers $n,m,w$, any type of instruction kinds with a designated
default, any function assigning to each kind a Boolean-list function, any program of
instructions, and any indices $i,j$, the circuit built from the program with output
indices $i,j$ is well formed, has no repeated input to any gate, uses exactly one input
for each NOT gate, and every gate has at most two inputs.
\end{theorem}

\begin{theorem}[Size of the program circuit]
\lean{BoolCircuit.progCircuit_size}
\label{BoolCircuit.progCircuit_size}
\leanok
\uses{BoolCircuit.DAGCircuit.size, BoolCircuit.GInstr, BoolCircuit.length_progGates, BoolCircuit.progBase, BoolCircuit.progCircuit}
\difficulty{3}
For any natural numbers $n,m,w$, any type of instruction kinds with a designated
default, any function assigning to each kind a Boolean-list function, any program of
instructions, and any indices $i,j$, the size of the circuit built from the program with
output indices $i,j$ equals $n+2$ plus the sum, over the instructions in the program, of
the widths of the gadgets associated to their kinds.
\end{theorem}

\begin{theorem}[Program circuit size bound]
\lean{BoolCircuit.progCircuit_size_le}
\label{BoolCircuit.progCircuit_size_le}
\leanok
\uses{BoolCircuit.DAGCircuit.size, BoolCircuit.GInstr, BoolCircuit.kindWidth, BoolCircuit.progBase, BoolCircuit.progCircuit, BoolCircuit.progCircuit_size}
\difficulty{4}
For any natural numbers $n,m,w$, any finite type of instruction kinds, any function
assigning to each kind a Boolean-list function, any program of instructions, and any
indices $i,j$, the size of the circuit built from the program with output indices $i,j$
is at most $n+2$ plus the program length multiplied by the maximum gadget width among
all instruction kinds.
\end{theorem}
